% Generated by Listas/make_statements.py from pratica-lecture-05.md, item 25. Do not edit.
\textbf{Source.} Mas-Colell, Whinston \& Green (1995), \emph{Microeconomic Theory}, Ch. 12 Exercises, Exercise 12.AA.2\^{}C, p.~435 {[}PDF 228{]}. Answer in the source: in the MWG solutions manual. Reproduced verbatim. \textbf{Added (not in the exercise):} the statement of Proposition 12.AA.4 and the definition of the minimax payoff it uses, both reproduced verbatim from p.~421, so the item can be worked without the book.



\begin{quote}

\textbf{12.AA.2\^{}C} Prove Proposition 12.AA.4. {[}\emph{Hint:} Consider a strategy profile of the following form: the players are to play an outcome path involving some pair \((q_1, q_2)\) in period 1 and \((q_1^*, q_2^*)\) in every period thereafter. If either player deviates, this outcome path is restarted.{]}



\emph{Added from p.~421.} To consider this issue, it is useful to let \(\underline{\pi}_i = \operatorname{Min}_{q_j}\left[\operatorname{Max}_{q_i} \pi_i(q_i, q_j)\right]\) denote player \(i\)'s \emph{minimax payoff}.



\textbf{Proposition 12.AA.4:} Consider an infinitely repeated game with \(\delta > 0\) and \(S_i \subset \mathbb{R}\) for \(i = 1, 2\). Suppose also that \(\pi_i(q)\) is differentiable at \(q^* = (q_1^*, q_2^*)\), with \(\partial \pi_i(q_i^*, q_j^*)/\partial q_j \neq 0\) for \(j \neq i\) and \(i = 1, 2\), and that \(\pi_i(q_1^*, q_2^*) > \underline{\pi}_i\) for \(i = 1, 2\). Then there is some SPNE with discounted payoffs to the two players of \((v_1', v_2')\) such that \((1 - \delta)v_i' < \pi_i(q_1^*, q_2^*)\) for \(i = 1, 2\).

\end{quote}



\textbf{Optional extras.}

- Gibbons 2.12 (p.~135): strategy, subgame and SPE in a repeated game; it is item 1 of the Lecture 4 list.

- Muñoz-Garcia Ch. 6 Ex. 7 (p.~299): the game of item 16 with a full worked answer --- check item 16 against it.

- Tadelis 10.12 (p.~219): the folk theorem on the trust game, after Lecture 6 introduces that game.

- Cabral 7.15 (p.~182): grim strategies on an advertising agreement, as a bridge to Lecture 6's collusion.



\begin{center}\rule{0.5\linewidth}{0.5pt}\end{center}

